% 第2章 2.5.4：条件期望作为分块平均（正交投影）。
% Ω=[0,1)，均匀概率，Z(ω)=ω^2。G_4 由四个等长区间生成：块平均 ((k+1)^3-k^3)/48，即 1/48,7/48,19/48,37/48；
% G_2 由两半生成：块平均 1/12, 7/12，恰为 G_4 块平均两两再平均（塔式性质）。
\begin{tikzpicture}[font=\figurefont\footnotesize]
\begin{axis}[width=76mm,height=42mm,axis lines=left,
  x axis line style={draw=FigureStroke,line width=.6pt},y axis line style={draw=FigureStroke,line width=.6pt},
  tick style={draw=FigureStroke},tick label style={font=\figurefont\scriptsize,text=FigureInk},
  xlabel={$\omega$},xlabel style={font=\figurefont\scriptsize,at={(axis description cs:1,0)},anchor=north west,yshift=1mm},
  xmin=0,xmax=1.02,ymin=0,ymax=1.05,xtick={0,0.25,0.5,0.75,1},ytick={0,0.25,0.5,0.75,1},
  ymajorgrids,grid style={draw=FigureGrid,line width=.4pt},
  legend style={draw=none,fill=none,font=\figurefont\scriptsize,cells={anchor=west},at={(1.02,0.5)},anchor=west}]
\addplot[draw=FigureInk,line width=1.1pt,domain=0:1,samples=80] {x^2};
\addplot[draw=FigureBlue,line width=1.3pt,const plot] coordinates {(0,0.020833) (0.25,0.145833) (0.5,0.395833) (0.75,0.770833) (1,0.770833)};
\addplot[draw=FigureGreen,line width=1.1pt,dashed,const plot] coordinates {(0,0.083333) (0.5,0.583333) (1,0.583333)};
\legend{$Z(\omega)=\omega^2$,$\mathbb E[Z\mid\mathcal G_4]$,$\mathbb E[Z\mid\mathcal G_2]$}
\end{axis}
\end{tikzpicture}
